\section{Arithmetic realization in split cyclotomic coefficient algebras}
\label{sec:prime-degree-local-cm}

Let \(p\ge5\) be prime, put \(g=(p-1)/2\), and consider
\[
 C_p:\ t^2-t=w^p,\qquad J_p=\operatorname{Jac}(C_p),
 \qquad F_p=\mathbf Q(\zeta_p).
\]
Write
\[
 \Gamma_p=(\mathbf Z/p)^\times,\qquad
 [a](\zeta_p)=\zeta_p^a,\qquad
 \mathcal A_p=O_{F_p}\otimes\mathbf Z_2,\qquad
 \mathcal U_p=1+2\mathcal A_p.
\]
Choose a primitive residue \(\gamma\bmod p\), put
\(s=[\gamma]\), and let \(c=[-1]=s^g\).
The groups relevant to the coefficient-two action are
\[
 \mathcal T_p=\{u\in\mathcal U_p:u\,c(u)=1\},
 \qquad \mathcal P_p=(s-1)\mathcal T_p.
\]
For an integral group-ring element \(A=\sum n_a[a]\),
the notation \(A(u)\) means \(\prod[a](u)^{n_a}\).
On pro-two unit groups this extends continuously to
coefficients in \(\mathbf Z_2\).

Let \(\rho_{p,2}\) be the established full integral
compact-homology representation of the cyclic cover with
monodromies \(1,1,-2\). Its scalar restriction to
\(M=\ker\lambda\) is denoted by \(\mathfrak h_p\).
The full duplication identity gives
\(\mathfrak h_p(N)\subset\mathcal T_p\).
We use the established full companion equality and
torus coinvariant identification in the form
\begin{equation}\label{eq:pdl-upper-input}
 \mathfrak h_p(R_{\rm quad})\subset\mathcal P_p.
\end{equation}
Explicitly, the two torus characters are the chosen
\(\nu_p\) and \(\psi_p\); their common kernel in
\(\mathcal T_p\) is \((s-1)\mathcal T_p\).
The arithmetic realization below is independent of
any assumption that two is primitive modulo \(p\).

\subsection{CM type and all local factors}

\begin{lemma}\label{lem:pdl-reflex-and-local}
The CM type of \(J_p\) is \(\{[1],\ldots,[g]\}\).
It is primitive, its reflex field is \(F_p\), and its
reflex norm on coefficient-two units is the group-ring map
\[
                           \Phi_p^*=\sum_{a=1}^{g}[a^{-1}].
\]
For every prime \(v\mid2\) of \(F_p\), let
\(\iota_v:O_{F_p,v}^\times\to\mathcal A_p^\times\)
put a unit in the \(v\)-component and one in every other
component. A corresponding inertia subgroup \(I_v\subset H\)
satisfies the exact arithmetic identity
\begin{equation}\label{eq:pdl-one-inertia}
                \rho_{J_p,2}(I_v)
                    =\Phi_p^*(\iota_v O_{F_p,v}^\times).
\end{equation}
Consequently, products of elements of these conjugate
inertia subgroups realize \(\Phi_p^*(u)\) for every
\(u\in\mathcal U_p\).
\end{lemma}

\begin{proof}
The regular differentials
\[
                   \frac{w^{j-1}\,dw}{2t-1},\qquad 1\le j\le g,
\]
have deck eigenvalues \(\zeta_p^j\). The integral deck
action embeds \(O_{F_p}\) in \(\operatorname{End}_{F_p}(J_p)\),
and its degree is \(p-1=2\dim J_p\). This proves the CM
assertion and the displayed type.
Its stabilizer in \(\Gamma_p\) is trivial: the sum of the
residues in the type is
\[
             \sum_{j=1}^{g}j=\frac{p^2-1}{8}
                                      \equiv-\frac18\pmod p,
\]
which is nonzero. If multiplication by \(a\) preserved
the type, comparison of sums would force \(a=1\).
The reflex field is therefore \(F_p\). An induced type
from a proper CM subfield would have a nontrivial stabilizer,
so the type is primitive. The inverse-embedding reflex norm
is precisely \(\Phi_p^*\), by
\cite[Example 1.19, p.~14; Proposition 1.26, p.~17]{MilneCM2020}.

The curve has good reduction at two. Besides its displayed
affine chart, use \(u=w^{-1}\), \(v=t/w^{g+1}\), giving
\[
                          v^2-u^{g+1}v=u.
\]
The two monic quadratic charts glue to a finite flat cover
of \(\mathbf P^1_{\mathbf Z_2}\). The first has derivative
\(2t-1\), a unit modulo two. In the second chart, vanishing
of the derivative in \(v\) modulo two forces \(u=0\),
where the derivative in \(u\) is \(-1\). Thus the proper
model is smooth, and its Jacobian has good reduction.

The field \(F_p\) is unramified at two. Each \(I_v\)
is consequently an inertia group of \(\mathbf Q\) at two
contained in \(G_{F_p}\). Local and global reciprocity
identify its abelian image with the global Artin image
of ideles supported on \(O_{F_p,v}^\times\).
For such an idele \(s_v\) and a corresponding \(b_v\in I_v\),
Milne's main theorem
\cite[Theorem 9.10, p.~78]{MilneCM2020} gives
\[
                   \rho_{J_p,f}(b_v)=A_v N_\Phi(s_v),
\]
where \(A_v:J_p\to b_vJ_p\) is an \(F_p\)-linear
quasi-isogeny. Here \(b_vJ_p=J_p\), since the Jacobian
is defined over \(\mathbf Q\), and \(b_v\) fixes the CM
field. At every auxiliary odd Tate coefficient, good
reduction at two makes the left side one, while the
reflex norm of \(s_v\) also has component one.
Faithfulness of the Tate action on quasi-endomorphisms
therefore gives \(A_v=1\). This proves
\eqref{eq:pdl-one-inertia}, on all units and without a
finite local correction. We use Milne's Artin convention
from printed p.~76; its reversal inverts the parameter
and does not change any image.

Every \(u\in\mathcal U_p\) is a product of the separate
\(\iota_v(u_v)\). Choose a corresponding \(b_v\) for each
factor and take their product \(b(u)\). The factors fix
\(F_p\), so their Tate images commute and multiply to
\(\Phi_p^*(u)\). No assertion that one local inertia
group supplies all factors is used.
\end{proof}

Here is an explicit description of the local product.
Let \(d=\operatorname{ord}_p(2)\) and \(r=(p-1)/d\).
Then \(\mathcal A_p\) is a product of \(r\) copies of
the unramified degree-\(d\) integer ring \(O_d\).
Choose representatives \(e_i\) for
\(\Gamma_p/\langle2\rangle\), and let \(\sigma\)
be Frobenius on \(O_d\). In these coordinates,
\begin{equation}\label{eq:pdl-frobenius-products}
 (\Phi_p^*(u))_i
    =\prod_{a=1}^{g}\sigma^{k(i,a)}(u_{j(i,a)}),
 \qquad e_i a^{-1}=2^{k(i,a)}e_{j(i,a)}\pmod p.
\end{equation}
This formula keeps every Frobenius and permutation factor.
If \(d\) is even, complex conjugation acts in each
component as \(\sigma^{d/2}\); \(\mathcal T_p\) is a
product of \(r\) relative norm-one principal-unit groups.
If \(d\) is odd, conjugation pairs components, and
\(\mathcal T_p\simeq(1+2O_d)^{r/2}\).

\subsection{An exact inertia-generated arithmetic subgroup}

\begin{proposition}\label{prop:pdl-arithmetic-realization}
For every \(u\in\mathcal U_p\), there is an element
\(w_p(u)\in V_{\rm quad}\cap H\) such that
\begin{equation}\label{eq:pdl-witness-values}
 \mathfrak h_p(w_p(u))=(s-1)^2\Phi_p^*(u),\qquad
 \mathfrak h_q(w_p(u))=1
        \quad\text{for every odd prime }q\ne p.
\end{equation}
Every original \(\chi\in X_2\) vanishes on this
entire family. In particular the image
\((s-1)^2\Phi_p^*(\mathcal U_p)\) is realized
arithmetically with all other prime-degree coordinates zero.
\end{proposition}

\begin{proof}
Choose \(a\in I_p\) with \(\lambda(a)\bmod p=\gamma\).
It has \(\lambda_2(a)=1\). Obtain \(b(u)\) as in
Lemma~\ref{lem:pdl-reflex-and-local}, and put
\[
                  z=[a,b(u)],\qquad z_2=z^{e_2},
                         \qquad w_p(u)=[a,z_2].
\]
The idempotent exponent \(e_2\in\widehat{\mathbf Z}\)
is one at two and zero elsewhere; powers are taken in
the procyclic closure of the element. Both \(z\) and
\(z_2\) have full multiplier one.

For the arithmetic Jacobian, conjugation by \(a\) acts
on the coefficient algebra by \(s\). Consequently
\[
                    \rho_{J_p,2}(z)=(s-1)\Phi_p^*(u).
\]
This value is pro-two because \(u\in\mathcal U_p\).
The based-cover action and the Jacobian action may differ
by a deck \(C_p\) factor. Indeed two lifts of the same
outer action differ by conjugation in the base geometric
group; modulo an inner automorphism of its covering
subgroup, this is exactly a deck transformation. Inner
automorphisms have trivial action on compact homology.
The discrepancy is a crossed map to \(\mu_p\), and
on the cyclotomic kernel it is a homomorphism.
It therefore vanishes on the pro-two element \(z_2\).
Thus
\[
                    \mathfrak h_p(z_2)=(s-1)\Phi_p^*(u).
\]
The full cover representation now computes its second
commutator by \(s-1\), proving the first equality in
\eqref{eq:pdl-witness-values}.

For membership in \(N\), use the coefficient-two Kummer
cocycle of two. It vanishes on \(a\in I_p\), since
\(p\ne2\). The cocycle commutator formula
\[
 \rho_q^I([h,k])=(1-\lambda_k)\rho_q^I(h)
                              +(\lambda_h-1)\rho_q^I(k)
\]
therefore gives \(\rho_{2,2}^I(z)=0\), because
\(\lambda_2(a)=1\). Taking \(e_2\) kills all other
coefficient components, so \(z_2\in N\cap H\).
Every rational Kummer coordinate of \(z_2\) is supported
at two. Hence
\[
                \rho_q^I(w_p(u))
                   =(\lambda(a)-1)\rho_q^I(z_2)=0
                           \quad\text{for every }q,
\]
and \(w_p(u)\in D\cap H\).

The based pro-\(\ell\) action of an inertia group
at two is trivial for \(\ell\ne2\), and the same is
true of every conjugate of that inertia group.
Thus \(\Phi_p(b(u))=1\). For \(\ell\ne p\), one
has \(\Phi_\ell(a)=1\). This is the automorphism-level
unramifiedness assertion with the rational tangential
base point in \cite[\S8.5, pp.~151--152]{HainMatsumoto2003}.
Consequently every \(\Phi_\ell\) kills \(z,z_2,w_p(u)\).

Since \(p\ge5\), inertia at \(p\) fixes both
\(\mathbf Q(i)\) and \(\mathbf Q(\sqrt{-3})\).
The element \(z_2\) fixes these fields as well, because
it has full multiplier one. Their harmonic and cubic
elliptic Tate actions lie in commutative CM algebras,
so their Tate commutators are trivial. The cubic frame's
quartic Kummer correction of three is trivial on \(D\).
Hence both complete corrected frames of \(w_p(u)\)
are one. The cubic radical--CM defect \(d_{(3)}\)
vanishes on the arithmetic subgroup under consideration.
This proves all the word, radical, defect, and full
coefficient-two/odd-frame conditions.

For any \(\chi\in X_2\), conjugation invariance gives
\[
                  \chi(w_p(u))
                    =\chi(a z_2a^{-1})-\chi(z_2)=0.
\]
This kills every chosen quadratic coordinate, including
the whole odd-pair family. Thus \(w_p(u)\in V_{\rm quad}\cap H\).

Finally fix an odd prime \(q\ne p\). Both \(a\) and
\(z_2\) have cyclotomic value one modulo \(q\): for
\(a\) this follows because \(\mathbf Q(\zeta_q)\) is
unramified at \(p\). Their degree-\(q\) homology actions
therefore commute with the deck generator. Its integral
centralizer is the commutative algebra
\(O_{F_q}\otimes\mathbf Z_2\): the nontrivial deck
eigencharacters occur once, and compact homology is
free of rank one over this algebra. The two actions
commute, proving the second equality in
\eqref{eq:pdl-witness-values}. This is an identity of the
full based-cover representations; no further
arithmetic or tangential comparison is needed for it.
\end{proof}

\subsection{The split prime seven}

For \(p=7\), let \(O_3\) be unramified cubic over
\(\mathbf Z_2\), with Frobenius \(\sigma\). Use the two
embeddings represented by \(1,-1\). Then
\[
 \mathcal A_7=O_3\times O_3,\qquad
 c(a,b)=(b,a),\qquad
 s(a,b)=(\sigma^2 b,\sigma^2 a),\quad \gamma=3.
\]
In particular \(s^2=\sigma\) componentwise and \(s^3=c\).
The reflex norm is
\begin{equation}\label{eq:pdl-seven-reflex}
 \Phi_7^*(a,b)
       =(a\sigma^2(a)\sigma(b),\ b\sigma^2(b)\sigma(a)).
\end{equation}
For a unit in the first factor, this gives
\[
 (s-1)\Phi_7^*(u,1)
                         =((\sigma^2u)^{-1},\sigma^2u).
\]
As \(u\) ranges over \(1+2O_3\), these are every
element of \(\mathcal T_7\). This verifies the exact
split local image directly, including its first unit layer.

Identifying \(\mathcal T_7\) with \(1+2O_3\) by
\(t\mapsto(t,t^{-1})\), one has
\[
 \mathcal P_7=(1+\sigma^2)(1+2O_3)
  =\{t\in1+2O_3:N_{O_3/\mathbf Z_2}(t)
                                  \in(\mathbf Z_2^\times)^2\}.
\]
To verify the last equality, the norm of
\(t\sigma^2(t)\) is a square. Conversely suppose
\(N(t)=v^2\), and choose the unique rational unit \(a\)
with \(a^3=v\). Then \(t/a^2\) has norm one. On this
norm-one subgroup, \(1+\sigma^2=-\sigma\), an
automorphism; and \(a^2=(1+\sigma^2)(a)\).
This proves the equality on the entire principal-unit group.
The two existing torus characters here are
\(a(Nt)\) and \(b(Nt)\), where
\(a(v)=(v^2-1)/8\pmod2\) and \(b(v)=(v-1)/2\pmod2\).
Thus their common kernel is exactly \(\mathcal P_7\),
also verifying \eqref{eq:pdl-upper-input} directly in this case.
The group \(\mathcal P_7\) has rank three and no torsion:
the only principal-unit torsion is \(-1\), whose cubic
norm is \(-1\), not a square. Hence
\(\mathcal P_7\simeq\mathbf Z_2^3\).

It follows from Proposition~\ref{prop:pdl-arithmetic-realization}
that all four images on \(R_{\rm quad},V_{\rm quad}\)
and their arithmetic intersections are \(\mathcal P_7\).
The same is true after imposing the degree-five kernel,
because the degree-seven witnesses have degree-five value one.
Thus adjoining the degree-seven kernel after degree five
gives a further strict reduction, with exact quotient
\(\mathbf Z_2^3\), already arithmetically.

\subsection{Simultaneous consequence of the integral unit criterion}

We formulate the arithmetic consequence of the explicit
CM-type unit calculation. Suppose
\[
 (s-1)\Phi_p^*=(c-1)Q_p
       \quad\text{in }\mathbf Z[\Gamma_p],
\]
and the image of \(Q_p\) is a unit in
\(\mathbf Z_2[s]/(s^g+1)\). Product-algebra unit Hilbert
90 gives \((c-1)\mathcal U_p=\mathcal T_p\): in a
self-conjugate unramified factor use ordinary unit
Hilbert 90 and remove its invariant Teichmuller part;
in a conjugate pair use \((1,t)\).
The operator \(Q_p\) is an automorphism on
\(\mathcal T_p\), since \(c\) acts there by inversion.
It follows that
\begin{equation}\label{eq:pdl-unit-surjectivity}
                     (s-1)\Phi_p^*(\mathcal U_p)=\mathcal T_p.
\end{equation}
Combining \eqref{eq:pdl-unit-surjectivity},
\eqref{eq:pdl-upper-input}, and
Proposition~\ref{prop:pdl-arithmetic-realization} gives
all four exact images \(\mathcal P_p\).
For seven the group-ring identity is
\[
 \Phi_7^*=1+s^4+s^5,\qquad
 (s-1)\Phi_7^*=(s^3-1)(-s),
\]
so its unit criterion requires no numerical certificate.

\begin{theorem}\label{thm:simultaneous-cyclotomic-cm}
Let \(\Pi\) be any set of primes at least five satisfying the equivalent
unit conditions of Theorem~\ref{thm:cm-type-torus-unit}, and put
\[
 \mathcal R_\Pi=R_{\rm quad}\cap\bigcap_{p\in\Pi}\ker\mathfrak h_p,
 \qquad \mathcal V_\Pi=V_{\rm quad}\cap\mathcal R_\Pi.
\]
These subgroups are closed and normal in \(G\). The joint image of the
old subgroups, and of either arithmetic intersection, is exactly
\begin{equation}\label{eq:pdl-simultaneous}
       \prod_{p\in\Pi}\mathcal P_p.
\end{equation}
Both old-to-new quotients are this product, also on arithmetic
intersections. Moreover \(G=\mathcal R_\Pi H\); full \(B\) is preserved
on \(\mathcal R_\Pi\), and \(B_{\max2}\) is preserved on
\(\mathcal V_\Pi\). The exact original restriction kernel on either
subgroup is \(\mathcal S=s(Y_2)\). Its complete image is the invariant
quadratic characters killing that subgroup's arithmetic intersection,
canonically \(\operatorname{Hom}_{\rm cts}(B,\mathbf F_2)\).
No nonzero continuous quadratic character of the displayed product,
pulled back to \(R_{\rm quad}\), extends to a member of \(X_2\).
\end{theorem}

\begin{proof}
The normalizer acts on each torus through \(\Gamma_p\), preserving its
identity. Thus the two kernels are closed and normal. The upper bound
and the axis witnesses have just been proved.
For every finite set of coordinates, multiply the arithmetic
witnesses from Proposition~\ref{prop:pdl-arithmetic-realization}:
each has arbitrary prescribed value on its own coordinate
and value one on all the others. This gives every finite
projection of \eqref{eq:pdl-simultaneous}. The image of
the compact arithmetic subgroup is closed, so these finite
surjections give the whole infinite product.

Arithmetic correction now gives
\[
 G=\mathcal R_{\Pi}H,\qquad
 B\simeq \mathcal R_{\Pi}/
 \overline{[G,\mathcal R_{\Pi}](\mathcal R_{\Pi}\cap H)}.
\]
The quotient \(\mathcal R_{\Pi}/\mathcal V_{\Pi}\)
embeds in the pro-odd group \(R_{\rm quad}/V_{\rm quad}\).
The central coprime extension lemma therefore preserves
\(B_{\max2}\) on \(\mathcal V_{\Pi}\).
The original restriction kernel on either subgroup remains
\(\mathcal S\). Indeed an original character vanishing on the new
kernel factors through \eqref{eq:pdl-simultaneous} on the
old subgroup. It vanishes on every arithmetic axis witness.
The finite-support subgroup is dense in the product, so
continuity forces that factored character to be zero.
The complete restricted images are again precisely the
invariant quadratic characters killing the arithmetic
intersection, canonically \(\operatorname{Hom}(B,\mathbf F_2)\),
by the complete quadratic splitting, ordinary gluing on
\(\mathcal R_{\Pi}\), and pro-odd extension for
\(\mathcal V_{\Pi}\).

A nonzero continuous quadratic character of the product is nonzero on
some arithmetic axis combination, whereas every original \(X_2\)
character vanishes on all those witnesses. This proves the nonextension
assertion.
\end{proof}

\subsection{Two limitations of a field-only generalization}

First, one inertia factor need not fill the torus. For
\(p=31\), one has \(d=5\), \(r=6\), and
\(\operatorname{rank}_{\mathbf Z_2}\mathcal T_{31}=15\).
One inertia image is a quotient of the unit group of an
unramified degree-five local field and therefore has
pro-two rank at most five. Applying fixed group-ring
operators cannot increase that rank. Thus replacing all
local factors by one inertia group cannot yield the
rank-fifteen target. The product construction is essential.

Second, \(\mathcal P_p\) need not be torsion-free. Set
\[
 e_p=\begin{cases}r,&d\text{ even},\\r/2,&d\text{ odd}.\end{cases}
\]
The sign torsion of \(\mathcal T_p\) is \((C_2)^{e_p}\),
and \(s\) cycles these sign coordinates. Its difference
has image \((C_2)^{e_p-1}\). This is all torsion in
\(\mathcal P_p\): if \((s-1)t\) has torsion image,
then \((s-1)\log t=0\). The invariant vector \(\log t\)
is diagonal and rational, while \((1+c)\log t=0\),
so it is zero. Thus \(t\) itself is sign torsion.
The rational anti-conjugation space has dimension \(g\),
and \(s-1\) is invertible on it. Hence, noncanonically,
\[
                  \mathcal P_p\simeq
                   \mathbf Z_2^g\times(C_2)^{e_p-1}.
\]
For example, at \(p=17\), \(d=8\) and \(r=2\),
so the target has a nontrivial sign-torsion factor.
All simultaneous assertions above retain the exact group
\(\mathcal P_p\); they do not discard this torsion.
